\chapter{Chapter09 — KKL Counterexample}
%%% generated by the blueprint pipeline from FABL.Chapter09.KKLCounterexample
\section{Overview}

Declarations of the Lean module \texttt{FABL.Chapter09.KKLCounterexample}.

\section{Declarations}

\begin{definition}[Normalized Rademacher sum]
\lean{FABL.kklNormalizedRademacherSum}
\statementsource{Analysis-of-Boolean-Functions-by-Ryan-ODonnell}{
  pdf-ffbe9c7b3306-p119-b004,
  pdf-ffbe9c7b3306-p273-b007
}
\label{FABL.kklNormalizedRademacherSum}
\leanok
\uses{FABL.linearForm}
For $n\in\mathbb N$, the normalized Rademacher sum on the sign cube is the linear form
whose coefficient at each coordinate is $1/\sqrt n$.
\end{definition}

\begin{definition}[Truncation to the unit interval]
\lean{FABL.truncateUnit}
\statementsource{Analysis-of-Boolean-Functions-by-Ryan-ODonnell}{pdf-ffbe9c7b3306-p273-b007}
\label{FABL.truncateUnit}
\leanok
For each real number $z$, define its truncation to $[-1,1]$ to be its projection onto
that interval.
\end{definition}

\begin{definition}[Truncated normalized Rademacher sum]
\lean{FABL.truncatedNormalizedRademacherSum}
\statementsource{Analysis-of-Boolean-Functions-by-Ryan-ODonnell}{pdf-ffbe9c7b3306-p273-b007}
\label{FABL.truncatedNormalizedRademacherSum}
\leanok
\uses{FABL.kklNormalizedRademacherSum, FABL.truncateUnit}
For $n\in\mathbb N$, the truncated normalized Rademacher sum on the sign cube is the
normalized Rademacher sum followed by projection onto $[-1,1]$.
\end{definition}

\begin{theorem}[Second moment of the normalized sign sum]
\lean{FABL.normalizedRademacherSum_secondMoment}
\label{FABL.normalizedRademacherSum_secondMoment}
\leanok
\uses{FABL.expect_linearForm_sq, FABL.kklNormalizedRademacherSum}
\difficulty{3}
For every positive integer $n$, if $x$ is uniformly distributed on $\{−1,1\}^n$, then
the second moment of $n^{-1/2}\sum_{i=1}^n x_i$ is $1$.
\end{theorem}

\begin{theorem}[Fourth moment bound for the normalized sum]
\lean{FABL.normalizedRademacherSum_fourthMoment_le_three}
\label{FABL.normalizedRademacherSum_fourthMoment_le_three}
\leanok
\uses{FABL.expect_linearForm_fourth, FABL.kklNormalizedRademacherSum}
\difficulty{4}
Let $n$ be a positive natural number, and choose $x$ uniformly from $\{-1,1\}^n$. The
fourth moment of the normalized coordinate sum satisfies
\[\mathbb E\left[\left(\frac{1}{\sqrt n}\sum_{i=1}^n x_i\right)^4\right]\le 3.\]
\end{theorem}

\begin{definition}[Sign-cube negation]
\lean{FABL.negateSignCubeEquiv}
\label{FABL.negateSignCubeEquiv}
\leanok
For each natural number $n$, the map on sign vectors of length $n$ that sends each
coordinate $x_i$ to $-x_i$ is an equivalence from the sign cube to itself.
\end{definition}

\begin{theorem}[Oddness of unit truncation]
\lean{FABL.truncateUnit_neg}
\label{FABL.truncateUnit_neg}
\leanok
\uses{FABL.truncateUnit}
\difficulty{4}
For every real number $z$, projection onto the interval $[-1,1]$ is odd: truncating $-z$
to this interval gives the negative of truncating $z$.
\end{theorem}

\begin{theorem}[Sign symmetry of a truncated sum]
\lean{FABL.truncatedNormalizedRademacherSum_neg}
\label{FABL.truncatedNormalizedRademacherSum_neg}
\leanok
\uses{FABL.negateSignCubeEquiv, FABL.normalizedRademacherSum_neg, FABL.truncateUnit_neg, FABL.truncatedNormalizedRademacherSum}
\difficulty{1}
For every $n\ge 0$ and every sign vector $x\in\{\pm1\}^n$, let
$S_n(x)=n^{-1/2}\sum_{i=1}^n x_i$, with the coefficient interpreted as the real
reciprocal of $\sqrt n$. Truncating a real number to the interval $[-1,1]$ is odd, so
truncating this normalized sum after negating every coordinate gives the negative of its
truncated value at $x$.
\end{theorem}

\begin{theorem}[Mean of a truncated normalized sign sum]
\lean{FABL.mean_truncatedNormalizedRademacherSum}
\label{FABL.mean_truncatedNormalizedRademacherSum}
\leanok
\uses{FABL.mean, FABL.negateSignCubeEquiv, FABL.truncatedNormalizedRademacherSum, FABL.truncatedNormalizedRademacherSum_neg}
\difficulty{4}
For every nonnegative integer $n$, the mean on the sign cube of the truncation to
$[-1,1]$ of the normalized Rademacher sum with all $n$ coefficients equal to $1/\sqrt n$
is zero.
\end{theorem}

\begin{theorem}[Lower bound after unit truncation]
\lean{FABL.truncateUnit_sq_ge_half_of_sq_gt_half}
\label{FABL.truncateUnit_sq_ge_half_of_sq_gt_half}
\leanok
\uses{FABL.truncateUnit}
\difficulty{4}
For every real number $z$ whose square is strictly greater than $1/2$, the square of its
projection onto $[-1,1]$ is at least $1/2$.
\end{theorem}

\begin{theorem}[Monotonicity of event probability]
\lean{FABL.pmfEventProbability_mono}
\label{FABL.pmfEventProbability_mono}
\leanok
\uses{FABL.pmfEventProbability, FABL.pmfExpectation_mono}
\difficulty{4}
Let $\Omega$ be a finite set, let $p$ be a probability mass function on $\Omega$, and
let $P,Q\subseteq\Omega$ satisfy $P\subseteq Q$. Then the probability assigned by $p$ to
$P$ is at most the probability assigned to $Q$.
\end{theorem}

\begin{theorem}[Variance lower bound for a truncated Rademacher sum]
\lean{FABL.variance_truncatedNormalizedRademacherSum_ge}
\label{FABL.variance_truncatedNormalizedRademacherSum_ge}
\leanok
\uses{FABL.kklNormalizedRademacherSum, FABL.mean, FABL.mean_truncatedNormalizedRademacherSum, FABL.mul_pmfEventProbability_le_expectation, FABL.normalizedRademacherSum_fourthMoment_le_three, FABL.normalizedRademacherSum_secondMoment, FABL.pmfEventProbability, FABL.pmfEventProbability_mono, FABL.pmfEventProbability_nonneg, FABL.pmfExpectation, FABL.pmfExpectation_uniformPMF_eq_expect, FABL.pmfPaleyZygmund, FABL.truncateUnit_sq_ge_half_of_sq_gt_half, FABL.truncatedNormalizedRademacherSum, FABL.uniformPMF, FABL.variance}
\difficulty{5}
For every positive integer $n$, let $X$ be uniform on $\{-1,1\}^n$ and set
$S_n(X)=n^{-1/2}\sum_{i=1}^n X_i$. The variance of $S_n(X)$ after truncating its value
to the interval $[-1,1]$ is at least $1/24$.
\end{theorem}

\begin{theorem}[Bound on the truncated value]
\lean{FABL.abs_truncateUnit_le_one}
\label{FABL.abs_truncateUnit_le_one}
\leanok
\uses{FABL.truncateUnit}
\difficulty{3}
For every real number $z$, its projection onto the interval $[-1,1]$ has absolute value
at most $1$.
\end{theorem}

\begin{theorem}[Effect of changing one coordinate]
\lean{FABL.normalizedRademacherSum_setCoordinate_sub}
\label{FABL.normalizedRademacherSum_setCoordinate_sub}
\leanok
\uses{FABL.kklNormalizedRademacherSum, FABL.linearForm, FABL.setCoordinate, FABL.signValue}
\difficulty{4}
For every natural number $n$, coordinate $i\in\{1,\ldots,n\}$, and sign vector
$x\in\{−1,1\}^n$, the normalized sum with coordinate $i$ set to $1$ minus the normalized
sum with that coordinate set to $-1$ equals $2/\sqrt n$.
\end{theorem}

\begin{theorem}[Derivative bound for a truncated Rademacher sum]
\lean{FABL.abs_discreteDerivative_truncatedNormalizedRademacherSum_le}
\label{FABL.abs_discreteDerivative_truncatedNormalizedRademacherSum_le}
\leanok
\uses{FABL.discreteDerivative, FABL.discreteDerivative_apply, FABL.kklNormalizedRademacherSum, FABL.normalizedRademacherSum_setCoordinate_sub, FABL.setCoordinate, FABL.truncateUnit, FABL.truncatedNormalizedRademacherSum}
\difficulty{4}
Let $n$ be a positive integer, let $i$ be a coordinate in the $n$-dimensional sign cube,
and let $x$ be a point of that cube. The absolute value of the discrete derivative in
coordinate $i$ of the truncated normalized Rademacher sum, evaluated at $x$, is at most
$1/\sqrt{n}$.
\end{theorem}

\begin{theorem}[Influence bound for a truncated normalized sum]
\lean{FABL.influence_truncatedNormalizedRademacherSum_le}
\label{FABL.influence_truncatedNormalizedRademacherSum_le}
\leanok
\uses{FABL.abs_discreteDerivative_truncatedNormalizedRademacherSum_le, FABL.discreteDerivative, FABL.influence, FABL.truncatedNormalizedRademacherSum}
\difficulty{4}
Let $n$ be a positive integer and $i$ a coordinate in $\{0,\ldots,n-1\}$. On the sign
cube $\{−1,1\}^n$, define $g(x)$ by truncating $\sum_{j=0}^{n-1}x_j/\sqrt n$ to the
interval $[-1,1]$. The influence of $g$ at $i$, the uniform average of its squared
discrete derivative, is at most $1/n$.
\end{theorem}

\begin{theorem}[Truncated normalized rademacher sum violates boolean kkl scale]
\lean{FABL.truncatedNormalizedRademacherSum_violates_boolean_kkl_scale}
\statementsource{Analysis-of-Boolean-Functions-by-Ryan-ODonnell}{
  pdf-ffbe9c7b3306-p097-b006,
  pdf-ffbe9c7b3306-p273-b007
}
\label{FABL.truncatedNormalizedRademacherSum_violates_boolean_kkl_scale}
\leanok
\uses{FABL.influence, FABL.influence_truncatedNormalizedRademacherSum_le, FABL.truncatedNormalizedRademacherSum, FABL.variance, FABL.variance_truncatedNormalizedRademacherSum_ge}
\difficulty{4}
% dataset-pipeline: stub -- prose pending build_informal_statements.py
Statement pending regeneration.
\end{theorem}
